\documentclass[11pt,a4paper]{article}

\usepackage[T1]{fontenc}
\usepackage[utf8]{inputenc}
\usepackage[brazilian]{babel}
\usepackage{amsmath, amssymb}
\usepackage{booktabs}
\usepackage{graphicx}
\usepackage{float}
\usepackage[margin=2.5cm]{geometry}
\usepackage{microtype}
\usepackage{xcolor}
\usepackage[hidelinks]{hyperref}
\usepackage{setspace}
\onehalfspacing
\graphicspath{{figuras_sem_motos/}}

\newcommand{\reais}[1]{R\$\,#1}

\title{\textbf{Valor monetário do patrimônio recuperado pela Força Municipal}\\[0.3em]
\large Um modelo probabilístico com distribuição de modelos e preços de mercado\\(Rio de Janeiro, março--junho de 2026)}
\author{Prefeitura do Rio --- Secretaria Municipal de Desenvolvimento Econômico \& COMPSTAT RIO}
\date{\today}

\begin{document}
\maketitle

\begin{abstract}
\noindent O patrimônio recuperado pela Força Municipal (SEGUR/Prefeitura do Rio) entre 15 de março e
8 de junho de 2026 --- 118 celulares, 19 bicicletas e 15 cordões --- é avaliado tratando o valor de
cada item como uma \emph{variável aleatória}: sorteia-se o modelo segundo a distribuição real do que
é roubado na cidade (ponderada por frequência) e o preço a partir de uma amostra de preços de
mercado, em duas bases --- \textbf{revenda} (usado) e \textbf{reposição} (novo). Por simulação de
Monte Carlo, o valor esperado do conjunto é de \textbf{\reais{169} mil} a preços de revenda (banda
de 90\%: \reais{149}--\reais{190} mil) e \textbf{\reais{275} mil} a preços de reposição
(\reais{242}--\reais{310} mil); o envelope extremo vai de \reais{12} mil a \reais{1{,}36} milhão.
Os celulares concentram cerca de 85\% do valor. Todas as cifras têm fonte pública (coleta em
16/06/2026).
\end{abstract}

% =====================================================================
\section{Introdução}
% =====================================================================
Entre 15 de março e 8 de junho de 2026, a Força Municipal (Secretaria Especial de Segurança
Urbana --- SEGUR, Prefeitura do Rio de Janeiro) reportou a recuperação de 118 celulares, 19
bicicletas e 15 cordões. Quanto vale, em reais, esse patrimônio?

A resposta não é um número único. Um celular recuperado pode ser um aparelho de entrada usado de
algumas centenas de reais ou um \emph{flagship} novo de mais de dez mil; um cordão pode ser uma
semijoia folheada ou uma peça de ouro. O valor de cada item recuperado é, portanto, uma
\emph{distribuição} de quantias plausíveis --- a probabilidade de que ele valha cada montante ---
e não apenas uma média.

Esse valor é uma variável aleatória, cuja distribuição se constrói a partir (i) da composição real
do que é roubado na cidade e (ii) de preços de mercado coletados por modelo, em duas bases ---
revenda (usado) e reposição (novo). Por \emph{revenda} entende-se a compra do bem \emph{usado}
(quanto ele valeria se revendido); por \emph{reposição}, a compra do equivalente \emph{novo} (o
custo de repô-lo). Cada categoria tem, assim, uma distribuição de probabilidade de valor, e o
conjunto, uma estimativa com intervalo de confiança --- tudo referenciado a fontes públicas.

% =====================================================================
\section{Metodologia}
% =====================================================================

\subsection{Modelo probabilístico}
Seja $i$ uma categoria de item ($\text{cel}$, $\text{bike}$, $\text{cord}$) e
$b\in\{\text{rev},\text{rep}\}$ a base de valor (revenda ou reposição). O valor de uma unidade
recuperada da categoria $i$ é a variável aleatória de \textbf{mistura}
\begin{equation}
  V_i^{\,b}\;\sim\;\sum_{m\in\mathcal{M}_i} w_{i,m}\,\hat F_{i,m}^{\,b},
  \label{eq:mix}
\end{equation}
em que $\mathcal{M}_i$ é o conjunto de modelos representativos, $w_{i,m}$ é o peso do modelo $m$
($\sum_m w_{i,m}=1$) e $\hat F_{i,m}^{\,b}$ é a distribuição empírica dos preços coletados para o
modelo $m$ na base $b$. Sorteia-se um modelo segundo os pesos e, em seguida, um preço a partir das
observações daquele modelo. Seu valor esperado e variância são
\begin{equation}
  \mathbb{E}[V_i^{b}]=\sum_m w_{i,m}\,\bar p_{i,m}^{\,b},
  \qquad
  \mathrm{Var}(V_i^{b})=\sum_m w_{i,m}\Big(\sigma_{i,m}^{b\,2}+\big(\bar p_{i,m}^{b}-\mathbb{E}[V_i^b]\big)^2\Big).
  \label{eq:moments}
\end{equation}
O valor total recuperado é a soma sobre todas as unidades:
\begin{equation}
  V^{\,b}=\sum_{i}\sum_{k=1}^{N_i} V_{i,k}^{\,b},
  \qquad
  \mathbb{E}[V^{b}]=\sum_i N_i\,\mathbb{E}[V_i^{b}],
  \qquad
  \mathrm{Var}(V^{b})=\sum_i N_i\,\mathrm{Var}(V_i^{b}),
  \label{eq:total}
\end{equation}
com $N_i$ a quantidade recuperada (painel da SEGUR). Nenhum ``fator de depreciação'' arbitrário é
aplicado: a base de revenda já é preço de usado (o mercado precifica o desgaste) e a de reposição
é o preço de novo.

\subsection{Pesos por frequência e bases de preço}
O item recuperado é o que se rouba na rua, então o peso $w_{i,m}$ de cada modelo vem de uma
\emph{fonte de frequência} (ranking de roubo, participação de mercado, mix de anúncios), e não de
uma escolha. O preço de cada modelo é uma \textbf{amostra de valores observados} (16/06/2026;
$n=323$: 216 de usado e 107 de novo), em duas bases --- revenda (usado) e reposição (novo) ---, com
cada anúncio ou listagem contando como um preço. OLX e Mercado Livre, que bloqueiam acesso
automatizado, entraram apenas por trechos de busca; os tamanhos por modelo estão em
\texttt{dados\_distribuicao.json}. Celulares, bicicletas e cordões são detalhados a seguir.

\subsection{Agregação por Monte Carlo}\label{sec:mc}
A distribuição de $V_i^b$ e de $V^b$ é obtida por simulação: sorteiam-se $N_i$ unidades por
categoria segundo a Eq.~\eqref{eq:mix}, somam-se, e repete-se 50.000 vezes (semente fixa, para
reprodutibilidade). Reporta-se o valor esperado e a banda de 90\% (percentis 5 e 95).

\subsection{Celulares}
Os pesos por marca vêm da frequência de roubo (Anuário do Fórum Brasileiro de Segurança Pública:
Samsung 37\%, Apple 25\%, Motorola 23\%, Xiaomi 10\%), ajustados pelo perfil socioeconômico (TIC
Domicílios) em direção aos modelos de entrada e intermediários, que formam a maior parte do parque.
Nove modelos representam a categoria, do Galaxy A0x ao iPhone Pro Max.

Os preços de usado vêm de páginas de varejo de seminovos lidas diretamente (Trocafy e Trocafone), em
que cada listagem é um preço --- 36 listagens do iPhone 11, 36 do iPhone 13, mais de 30 do Galaxy
S23, além de Galaxy A14, Moto G54 e iPhone 13 Pro Max ---; os modelos legados entram por busca. Os
preços de novo vêm do varejo (Magazine Luiza, Samsung, Apple). A distribuição é fortemente
assimétrica à direita: a massa fica nos aparelhos de entrada (revenda de algumas centenas de reais)
e a cauda nos \emph{flagships} (reposição acima de \reais{6.000}).

\subsection{Bicicletas}
As bicicletas são avaliadas como bens de uso pessoal --- as efetivamente compradas e furtadas de
pessoas físicas. Os pesos das convencionais seguem o mix de usados do Rio (OLX-RJ: Caloi
$\approx$70\%, à frente de Oggi e GTS): Caloi de entrada aro 26, outras nacionais, passeio genérica,
MTB aro 29 e infantil. Os preços vêm de anúncios de usado do Rio (OLX-RJ) e de páginas de seminovos
lidas diretamente (Las Magrelas, Decathlon), e do varejo no novo (Magazine Luiza, Mormaii, Caloi). A
bicicleta de sistema compartilhado (Tembici) é custo de sistema e fica fora da soma.

À mistura soma-se a \textbf{bicicleta elétrica}, que pode custar várias vezes mais que uma
convencional --- de \reais{1.400} (entrada) a \reais{18.000} (E-MTB), com a urbana em torno de
\reais{5.900} (varejo: Caloi, Oggi, Decathlon, Sense, Two Dogs); no usado, os seminovos vão de
\reais{1.500} a cerca de \reais{12.800}. A e-bike entra com peso de \textbf{7\%} das bicicletas: o
Instituto de Segurança Pública do RJ não separa elétrica de convencional no registro de furto, e a
fração vem de aproximações --- as e-bikes são 5,4\% da produção nacional de bicicletas (Abraciclo,
2024) e cerca de 2\% das vendas (Aliança Bike), com sobre-representação no roubo por alto valor (o
furto de bicicleta no Rio subiu 59\% em 2025, e o aplicativo 190RJ da Polícia Militar cadastra o
chassi ``especialmente em modelos elétricos''). A distribuição da bicicleta passa, assim, a ter
corpo nas convencionais (\reais{80}--\reais{2.000}) e cauda longa nas e-bikes.

Como o share de e-bikes não tem recorte oficial, o resultado é testado variando a fração $f_e$ em
torno do caso-base ($\mathbb{E}[V_{\text{bike}}]=\reais{981}$ em revenda, Tabela~\ref{tab:itens}):
\begin{center}\small
\begin{tabular}{lccc}
\toprule
$f_e$ (fração de e-bikes) & 3\% & $\approx$7\% (base) & 14\% \\
\midrule
$\mathbb{E}[V_{\text{bike}}]$ (revenda) & \reais{802} & \reais{981} & \reais{1.293} \\
\bottomrule
\end{tabular}
\end{center}
Sem as motocicletas, a bicicleta responde por cerca de 11\% do patrimônio
(Tabela~\ref{tab:contrib}); ainda assim, o total varia pouco com $f_e$ dentro dessa faixa.

\subsection{Cordões}
O valor do cordão segue a forma canônica das distribuições de bens furtados: fortemente
assimétrica à direita, com a massa em peças de baixo valor e uma cauda longa e fina nas peças mais
caras. O levantamento do Home Office britânico sobre o valor de bens roubados ilustra essa forma
--- a média (\pounds481) supera de longe a mediana ($\approx$\pounds100) e os 2\% de furtos mais
valiosos concentram 46\% do valor total (SHAW \emph{et al.}, 2015) ---, da mesma família
\emph{lognormal} de cauda pesada empregada para perdas na literatura atuarial (FREES, 2018; KLUGMAN
\emph{et al.}). Sob essa forma, a cauda fica contida em \reais{3.000} (revenda) e \reais{5.000}
(reposição), o que explica o desvio superar a média na Tabela~\ref{tab:itens}.

A composição acompanha a distribuição de mercado: a semijoia folheada responde pela maior parte das
peças em circulação --- lidera o varejo em unidades (EUROMONITOR, 2025), com Limeira como segundo
maior polo mundial do setor, em lotes de dezenas a centenas de milhares de peças (ETULAIN
\emph{et al.}, 2021; IBGM) ---, enquanto as peças de ouro, menos frequentes, concentram o valor. O
perfil do que circula na rua reforça a predominância do folheado: parte expressiva da população
evita usar objetos de valor em público (IBGE, 2021: 42,8\%) ou retira joias antes de sair
(FBSP/Datafolha, 2026: 26,8\%). No caso-base, cerca de 30\% das peças são de ouro --- proporção
coerente com a valorização do metal, que vem ampliando o roubo dirigido a joias (os roubos de
aliança em São Paulo mais que dobraram entre 2023 e 2025).

A participação do ouro é o parâmetro mais influente dessa categoria; o resultado é testado variando
a fração $f$ em torno do caso-base ($\mathbb{E}[V_{\text{cord}}]=\reais{462}$ em revenda,
Tabela~\ref{tab:itens}):
\begin{center}\small
\begin{tabular}{lccc}
\toprule
$f$ (fração de ouro) & 15\% & $\approx$30\% (base) & 40\% \\
\midrule
$\mathbb{E}[V_{\text{cord}}]$ (revenda) & \reais{260} & \reais{462} & \reais{588} \\
Total $V^{\text{rev}}$ & \reais{166} mil & \reais{169} mil & \reais{170} mil \\
\bottomrule
\end{tabular}
\end{center}
O resultado agregado é robusto: como os cordões representam cerca de 4\% do patrimônio
(Tabela~\ref{tab:contrib}), o valor total varia perto de 3\% mesmo dobrando $f$. As cifras de preço
--- folheado de \reais{15} a \reais{250} e cordões de ouro 18k de \reais{280} a \reais{3.000}, à
cotação de $\approx$\reais{530}/g --- estão documentadas no inventário de fontes.


% =====================================================================
\section{Resultados}
% =====================================================================

\subsection{Distribuição de valor por item}
A Figura~\ref{fig:itens} mostra a distribuição de valor de uma unidade recuperada de cada
categoria, nas duas bases. As formas são informativas: o celular é fortemente assimétrico à
direita (massa em aparelhos de entrada, cauda nos \emph{flagships}); a bicicleta tem corpo nas
convencionais e cauda longa nas elétricas; o cordão concentra a massa no folheado de baixo valor,
com cauda longa e fina nas peças de ouro (média bem acima da mediana). A
Tabela~\ref{tab:itens} resume os momentos.

\begin{figure}[H]
  \centering
  \includegraphics[width=\linewidth]{dist_itens}
  \caption{Distribuição de probabilidade do valor de uma unidade recuperada, por categoria
  (azul: revenda/usado; laranja: reposição/novo). Histograma das amostras com densidade KDE.
  Os painéis de cordões e bicicletas usam escala logarítmica no eixo $x$ para exibir o corpo de
  baixo valor e a cauda cara (ouro; e-bikes) simultaneamente.}
  \label{fig:itens}
\end{figure}

\begin{table}[H]
  \centering\small
  \begin{tabular}{llrrr}
    \toprule
    Item & base & $\mathbb{E}[V_i]$ & desvio & faixa p5--p95 \\
    \midrule
    Celulares  & revenda   & 1.212 & 944   & 530--3.500 \\
               & reposição & 1.964 & 1.641 & 619--6.500 \\
    Bicicletas & revenda   & 981   & 1.447 & 120--3.100 \\
               & reposição & 1.513 & 1.964 & 410--4.490 \\
    Cordões    & revenda   & 462   & 804   & 20--2.400 \\
               & reposição & 936   & 1.486 & 40--5.000 \\
    \bottomrule
  \end{tabular}
  \caption{Momentos da distribuição de valor por unidade (R\$), por base. Fonte: amostras de
  preço de mercado; simulação de Monte Carlo.}
  \label{tab:itens}
\end{table}

\subsection{Valor total recuperado}
Agregando (Eq.~\eqref{eq:total}), o valor esperado do patrimônio recuperado é de
\textbf{\reais{169} mil} a preços de revenda e \textbf{\reais{275} mil} a preços de reposição. A
Figura~\ref{fig:total} mostra as duas distribuições com suas bandas de 90\%
(\reais{149}--\reais{190} mil e \reais{242}--\reais{310} mil). A Tabela~\ref{tab:contrib} detalha a
contribuição por categoria: os \textbf{celulares} respondem por cerca de \textbf{85\%} do total.

\begin{figure}[H]
  \centering
  \includegraphics[width=0.92\linewidth]{dist_total}
  \caption{Distribuição do valor total recuperado por simulação de Monte Carlo (50.000 iterações);
  linha = valor esperado, faixa sombreada = banda de 90\%.}
  \label{fig:total}
\end{figure}

\begin{table}[H]
  \centering\small
  \begin{tabular}{lrrrr}
    \toprule
    Item & $N_i$ & Contrib.\ revenda & Contrib.\ reposição & \% (revenda) \\
    \midrule
    Celulares  & 118 & 143.016 & 231.752 & 84,8\% \\
    Bicicletas & 19  & 18.639  & 28.747  & 11,1\% \\
    Cordões    & 15  & 6.930   & 14.040  & 4,1\% \\
    \midrule
    \textbf{Total} & — & \textbf{$\approx$168.585} & \textbf{$\approx$274.539} & 100\% \\
    \bottomrule
  \end{tabular}
  \caption{Contribuição por categoria ($N_i\cdot\mathbb{E}[V_i]$) e total (R\$).}
  \label{tab:contrib}
\end{table}

\noindent A banda de 90\% reflete a incerteza de \emph{amostragem} do total. A incerteza
\emph{estrutural} é capturada pela distância entre as bases de revenda e reposição
(\reais{169}--\reais{275} mil) e pelo \textbf{envelope extremo} de \reais{12} mil (se cada item
fosse o modelo usado mais barato) a \reais{1{,}36} milhão (se cada item fosse o modelo novo mais
caro).

% =====================================================================
\section{Conclusões}
% =====================================================================
O patrimônio recuperado pela Força Municipal no período vale, em valor esperado, entre \reais{169}
mil (revenda) e \reais{275} mil (reposição), com envelope de \reais{12} mil a \reais{1{,}36}
milhão. O resultado é dominado pelos celulares ($\approx$85\%). Cada item tem não só um número, mas
uma \emph{distribuição} de valores plausíveis, ancorada na composição real do que se rouba na cidade
e em preços de mercado verificáveis. Períodos futuros requerem apenas a atualização dos $N_i$
(quantidades recuperadas).

% =====================================================================
\section*{Limitações}
% =====================================================================
\begin{itemize}\setlength\itemsep{0.1em}
  \item \textbf{Distribuição é proxy:} o \emph{mix} exato dos itens recuperados não foi divulgado;
        os pesos vêm da frequência de roubo/vendas, e a maioria das cifras é nacional, não recorte
        do Rio.
  \item \textbf{Amostras de preço, não censo:} OLX e Mercado Livre bloqueiam coleta automatizada;
        usaram-se varejo de seminovos, anúncios e joalherias.
  \item \textbf{Bicicletas:} avaliadas como bens de uso pessoal (varejo Caloi/Mormaii e usados
        OLX-RJ); inclui-se a bicicleta elétrica, cujo peso no roubo é estimado por aproximações (o
        ISP-RJ não separa elétrica de convencional) e testado por sensibilidade. Bicicletas de
        sistema compartilhado não entram na conta.
\end{itemize}

% =====================================================================
\section*{Fontes}
% =====================================================================
\small
\noindent Coleta em 16/06/2026. Inventário (formato F-T.N) em \texttt{inventario\_fontes.md};
amostras de preço por modelo (com $n$ e URL) em \texttt{dados\_distribuicao.json}.
(FÓRUM BRAS.\ DE SEGURANÇA PÚBLICA. ``Marcas de celulares mais roubadas'', via Exame/InfoMoney, 2024) ·
(STATCOUNTER. ``Mobile Vendor Share Brazil'', 2026) ·
(CETIC.BR. ``TIC Domicílios'', 2024) ·
(TROCAFY/TROCAFONE. ``Seminovos por listagem'', 2026) ·
(MAGAZINE LUIZA; SAMSUNG; APPLE; CALOI. ``Preços de varejo (novo)'', 2026) ·
(OLX. ``Bicicletas e celulares usados, RJ'', 2024--2026) ·
(OURO RIO DE JANEIRO. ``Cotação do ouro 18k'', 2026) ·
(SHEKINAH; PRATA PURA. ``Cordões folheados (preço)'', 2026; DR JOIAS. ``Ouro 18k leve 1,8--8\,g'', 2026) ·
(SHAW, O.\ \emph{et al.} ``Crime and the value of stolen goods'', Home Office Research Report 81, 2015) ·
(FREES, E.\ (ed.). ``Loss Data Analytics'', International Actuarial Association, 2018; KLUGMAN \emph{et al.} ``Loss Models'') ·
(IBGE. ``Vitimização e acesso à justiça, 2009'', 2010; SENASP/CRISP-UFMG. ``Pesquisa Nacional de Vitimização'', Min.\ Justiça, 2013) ·
(IBGE. ``PNAD Vitimização 2021'', 2022; FBSP/DATAFOLHA. ``Mudança de hábitos por insegurança'', 2026; EUROMONITOR. ``Jewellery in Brazil'', 2025) ·
(ETULAIN, C.\ \emph{et al.} ``Produção de semijoias em Limeira'', UNICAMP, 2021; IBGM).

\end{document}
